% Part: lambda-calculus
% Chapter: church-rosser
% Section: parallel-beta-reduction

\documentclass[../../../include/open-logic-section]{subfiles}

\begin{document}

\olfileid{lam}{cr}{pb}

\olsection{సమాంతర $\beta$-తగ్గింపు}

ఇక్కడ \emph{సమాంతర $\beta$-తగ్గింపు} అనే భావాన్ని పరిచయం చేసి,
దానికి చర్చ్--రోసర్ లక్షణం ఉందని నిరూపిస్తాం.

\begin{defn}[సమాంతర $\beta$-తగ్గింపు, $\bredpar$] \ollabel{defn:bredpar}
  పదాలపై సమాంతర తగ్గింపును ($\bredpar$) ఆగమన పద్ధతిలో ఇలా
  నిర్వచిస్తాం:
  \begin{enumerate}
    \item \ollabel{defn:bredpar1} $x \bredpar x$.
    \item \ollabel{defn:bredpar2} $N \bredpar N'$ అయితే $\lambd[x][N] \bredpar
      \lambd[x][N']$.
      \sourcecorrection{OLTELAMCRPB-001}{మూలంలోని ఈ నియమం పూర్వాపేక్షలో $N \xrightarrow{\beta} N'$ అని ఉంది. కానీ సమాంతర తగ్గింపు ఆగమన నిర్వచనంలో, అలాగే కింది \olref{thm:refl}, \olref{lem:comp}, \olref{lem:cont} నిరూపణల్లో, అవసరమైనది $N \bredpar N'$. మూలంలో లేని ఆ సమాంతర పూర్వాపేక్షను స్పష్టంగా పునరుద్ధరించాం; స్థిర ఆంగ్ల మూలాన్ని మార్చలేదు.}
    \item \ollabel{defn:bredpar3} $P \bredpar P'$, $Q \bredpar Q'$ అయితే $PQ \bredpar
      P'Q'$.
    \item \ollabel{defn:bredpar4} $N \bredpar N'$, $Q \bredpar Q'$ అయితే
      $(\lambd[x][N])Q \bredpar \Subst{N'}{Q'}{x}$.
  \end{enumerate}
\end{defn}

సమాంతర $\beta$-తగ్గింపులో ఒక పదంలోని ఎన్ని రెడెక్స్‌లనైనా
ఒకే దశలో సంకోచించవచ్చు. అయితే అసలు పదంలో ఉన్న రెడెక్స్‌లనే
సంకోచించవచ్చు; ఆ దశలో కొత్తగా ఏర్పడిన $\beta$-రెడెక్స్‌లను అదే
దశలో సంకోచించలేం. అందుకే ఇది $\beta$-తగ్గింపుకన్నా భిన్నం.
ఉదాహరణకు, $(\lambd[f][fx])(\lambd[y][y])$ను సమాంతర
$\beta$-తగ్గింపులో అదే పదంగా ఉంచవచ్చు లేదా
$(\lambd[y][y])x$కు తగ్గించవచ్చు; ఒకే దశలో~$x$కు
తగ్గించలేం. ఎందుకంటే ఆ రెడెక్స్ మొదటి సమాంతర తగ్గింపు
తరువాతే ఏర్పడుతుంది. అయితే రెండవ సమాంతర $\beta$-తగ్గింపు దశలో
$x$ వస్తుంది; సాధారణ $\beta$-తగ్గింపులో కూడా పదం~$x$కు
తగ్గుతుంది.

\begin{thm}\ollabel{thm:refl}
  $M \bredpar M$.
\end{thm}

\begin{proof}
  సాధన.
\end{proof}

\begin{prob}
  \olref[lam][cr][pb]{thm:refl}ను నిరూపించండి.
\end{prob}

\begin{defn}[$\beta$-సంపూర్ణ వికాసం]\ollabel{defn:bcd}
  $M$ పదం యొక్క \emph{$\beta$-సంపూర్ణ వికాసం}~$\bcd{M}$ను
  ఆగమన పద్ధతిలో ఇలా నిర్వచిస్తాం:
  \begin{align}
    \bcd{x} &= x \ollabel{defn:bcd1} \\
    \bcd{(\lambd[x][N])} &= \lambd[x][\bcd{N}] \ollabel{defn:bcd2}\\
    \bcd{(PQ)} &= \bcd{P}\bcd{Q} && \text{$P$ ఒక $\lambd$-అమూర్తీకరణ కాకపోతే}
    \ollabel{defn:bcd3} \\
    \bcd{((\lambd[x][N])Q)} &= \Subst{\bcd{N}}{\bcd{Q}}{x} \ollabel{defn:bcd4}
  \end{align}
\end{defn}

పదం యొక్క $\beta$-సంపూర్ణ వికాసం, దాని పేరుకు తగినట్లే,
“సంపూర్ణ సమాంతర తగ్గింపు.” సమాంతర $\beta$-తగ్గింపులో
ఒక రెడెక్స్‌ను సంకోచించకుండా వదలవచ్చు; సంపూర్ణ వికాసంలో
మాత్రం అసలు పదంలోని $\beta$-రెడెక్స్‌లన్నిటినీ సంకోచించాలి.
కాబట్టి $(\lambd[f][fx])(\lambd[y][y])$ యొక్క
సంపూర్ణ వికాసం $(\lambd[y][y])x$; అసలు పదం కాదు.

\begin{editorial}
  ఈ నిర్వచనానికి ఒక సమస్య ఉంది: ($\lambd$-)పదాలపై
  ప్రమేయాలను పునరావృతంగా ఎలా నిర్వచించాలో ఇంకా పరిచయం
  చేయలేదు. దీన్ని తరువాత సరిచేస్తాం.
\end{editorial}

\begin{lem}\ollabel{lem:comp}
  $M \bredpar M'$, $R \bredpar R'$ అయితే,
  $\Subst{M}{R}{y} \bredpar \Subst{M'}{R'}{y}$.
\end{lem}

\begin{proof}
  $M \bredpar M'$ యొక్క \tetoken{వ్యుత్పత్తి}{!!{derivation}}పై
  ఆగమన పద్ధతిలో నిరూపిస్తాం.
  \begin{enumerate}
    \item చివరి దశ \olref{defn:bredpar1}: సాధన.
    \item చివరి దశ \olref{defn:bredpar2}: అప్పుడు $M$ అనేది
      $\lambd[x][N]$, $M'$ అనేది $\lambd[x][N']$;
      ఇక్కడ $N \bredpar N'$. మనం నిరూపించవలసింది
      $\Subst{(\lambd[x][N])}{R}{y} \bredpar
      \Subst{(\lambd[x][N'])}{R'}{y}$, అంటే
      $\lambd[x][\Subst{N}{R}{y}] \bredpar
      \lambd[x][\Subst{N'}{R'}{y}]$.
      \sourcecorrection{OLTELAMCRPB-002}{మూలంలోని ఈ విప్పిన సూత్రం కుడివైపు $\Subst{N'}{R}{y}$ అని రాస్తుంది. ఉపసిద్ధాంతపు లక్ష్యం, అదే వాక్యంలోని విప్పకముందు సూత్రం, ఆగమన పరికల్పన అన్నీ కుడివైపు $R'$ కావాలని నిర్ధారిస్తున్నాయి; ఇక్కడ ఆ సూచికను సరిచేశాం.}
      ఇది \olref{defn:bredpar2}, ఆగమన పరికల్పనల నుంచి
      వెంటనే వస్తుంది.
    \item చివరి దశ \olref{defn:bredpar3}: సాధన.
    \item చివరి దశ \olref{defn:bredpar4}: $M$ అనేది
      $(\lambd[x][N])Q$, $M'$ అనేది $\Subst{N'}{Q'}{x}$.
      మనం నిరూపించవలసింది
      $\Subst{((\lambd[x][N])Q)}{R}{y} \bredpar
      \Subst{\Subst{N'}{Q'}{x}}{R'}{y}$, అంటే
      $(\lambd[x][\Subst{N}{R}{y}])\Subst{Q}{R}{y} \bredpar
      \Subst{\Subst{N'}{R'}{y}}{\Subst{Q'}{R'}{y}}{x}$.
      మూలం దీన్ని \olref{defn:bredpar4}, ఆగమన పరికల్పనల
      ఫలితంగా పేర్కొంటుంది.
      \sourcecorrection{OLTELAMCRPB-003}{మూలం ఈ రెండు సూత్రాలను “అంటే” అని సమానార్థకంగా చూపి, చివరి నిర్ధారణకు అవసరమైన ప్రతిస్థాపనల నిర్వచితత్వం, బద్ధ చరం $x$కు తాజా ప్రతినిధి ఎంపిక, ప్రతిస్థాపనల మార్పిడి సమర్థనను ఇవ్వలేదు. మునుపటి పాక్షిక ప్రతిస్థాపన, ఆల్ఫా-వర్గ నిర్మాణాల్లో OLTELAMALP-005–006 ఖాళీలూ తెరిచే ఉన్నాయి. ఈ నిరూపణ దశను మూల వాదనగా మాత్రమే ఉంచాం; పూర్తి నిరూపణగా ధ్రువీకరించడం లేదు.}
  \end{enumerate}
\end{proof}

\begin{prob}
  \olref[lam][cr][pb]{lem:comp} నిరూపణను పూర్తి చేయండి.
\end{prob}

\begin{lem}\ollabel{lem:cont}
  $M \bredpar M'$ అయితే $M' \bredpar \bcd{M}$.
\end{lem}
\begin{proof}
  $M \bredpar M'$ యొక్క \tetoken{వ్యుత్పత్తి}{!!{derivation}}పై
  ఆగమన పద్ధతిలో నిరూపిస్తాం.
  \begin{enumerate}
    \item చివరి నియమం \olref{defn:bredpar1}: సాధన.
    \item చివరి నియమం \olref{defn:bredpar2}: $M$ అనేది
      $\lambd[x][N]$, $M'$ అనేది $\lambd[x][N']$;
      ఇక్కడ $N \bredpar N'$. మనం చూపవలసింది
      $\lambd[x][N'] \bredpar
      \bcd{(\lambd[x][N])}$; అంటే \olref{defn:bcd2}
      ప్రకారం $\lambd[x][N'] \bredpar
      \lambd[x][\bcd{N}]$. ఇది \olref{defn:bredpar2},
      ఆగమన పరికల్పనల నుంచి వస్తుంది.
    \item చివరి నియమం \olref{defn:bredpar3}: $M$ అనేది
      $PQ$, $M'$ అనేది $P'Q'$; ఇక్కడ $P$, $Q$, $P'$, $Q'$
      కొన్ని పదాలు, $P \bredpar P'$, $Q \bredpar Q'$.
      ఆగమన పరికల్పన ప్రకారం $P' \bredpar \bcd{P}$,
      $Q' \bredpar \bcd{Q}$.
      \begin{enumerate}
        \item $P$ అనేది ఏవో $x$, $N$లకు $\lambd[x][N]$
          అయితే, $P'$ ఏదో $N'$కు $\lambd[x][N']$ అవుతుంది;
          ఇక్కడ $N \bredpar N'$. ఆగమన పరికల్పన ప్రకారం
          $N' \bredpar \bcd{N}$, $Q' \bredpar \bcd{Q}$.
          అప్పుడు \olref{defn:bredpar4} ప్రకారం
          $(\lambd[x][N'])Q' \bredpar
          \Subst{\bcd{N}}{\bcd{Q}}{x}$.
        \item $P$ ఒక $\lambd$-అమూర్తీకరణ కాకపోతే,
          \olref{defn:bredpar3} ప్రకారం $P'Q' \bredpar
          \bcd{P}\bcd{Q}$; \olref{defn:bcd3} ప్రకారం
          కుడివైపు పదమే $\bcd{PQ}$.
      \end{enumerate}
    \item చివరి నియమం \olref{defn:bredpar4}: $M$ అనేది
      $(\lambd[x][N])Q$; $M'$ అనేది $\Subst{N'}{Q'}{x}$.
      ఇక్కడ $x$, $N$, $Q$, $N'$, $Q'$ తగిన చరం,
      పదాలు; $N \bredpar N'$, $Q \bredpar Q'$.
      ఆగమన పరికల్పన ప్రకారం $N' \bredpar \bcd{N}$,
      $Q' \bredpar \bcd{Q}$. \olref{lem:comp} ప్రకారం
      $\Subst{N'}{Q'}{x} \bredpar
      \Subst{\bcd{N}}{\bcd{Q}}{x}$; కుడివైపు పదం
      సరిగ్గా $\bcd{((\lambd[x][N])Q)}$.
  \end{enumerate}
\end{proof}

\begin{prob}
  \olref[lam][cr][pb]{lem:cont} నిరూపణను పూర్తి చేయండి.
\end{prob}

\begin{thm}\ollabel{thm:cr}
  $\bredpar$కు చర్చ్--రోసర్ లక్షణం ఉంది.
\end{thm}

\begin{proof}
  \olref{lem:cont} నుంచి నేరుగా వస్తుంది.
\end{proof}

\end{document}
